%% kapa: exam=1 round=26 year=2015 period=2 form=A subject=appraisal_law subject_ko=부동산관계법규 seq=36 number=36
%% source: Q-Net 감정평가사 기출문제 (cst00309) artlSeq=5138001 file=extracted/1008/5138001/제26회_감정평가사 1차_기출/제26회_감정평가사 1차_기출/2교시_감평_A형.pdf
%% conversion: kordoc 4.24.51 -> tools/gen1.py (auto); review=proofread-sonnet; flags=none
\begin{question}{36}
측량ㆍ수로조사 및 지적에 관한 법령상 지목의 구분 및 설정방법 등에 관한 설명으로 옳지 않은 것은?
\begin{choices}
  \item 1필지가 둘 이상의 용도로 활용되는 경우에는 주된 용도에 따라 지목을 설정한다.
  \item 토지가 일시적 또는 임시적인 용도로 사용될 때에는 지목을 변경하지 아니한다.
  \item 원상회복을 조건으로 돌을 캐내는 곳 또는 흙을 파내는 곳으로 허가된 토지는 ‘잡종지’로 한다.
  \item 자연의 유수가 있거나 있을 것으로 예상되는 토지는 ‘하천’으로 한다.
  \item 육상에 인공으로 조성된 수산생물의 번식 또는 양식을 위한 시설을 갖춘 부지와 이에 접속된 부속 시설물의 부지는 ‘양어장’으로 한다.
\end{choices}
\end{question}
\answer{3}
